\chapter{Limitations and Threats to Validity}
\label{ch:limitations}

This chapter details the throughput limitations of the Kronos search core and outlines potential threats to the validity of our experimental findings, alongside mitigation steps.

\section{System Limitations}
The primary limitation of Kronos is its search throughput, which is constrained by the mailbox board representation. Unlike native C/C++ engines that utilize 64-bit integer bitboards updated via bitwise CPU operations, Kronos executes searches using object-based wrappers. This representation, combined with dynamic runtime overhead, limits throughput to approximately 5,000 NPS (Table~\ref{tab:cross_depth_elo}).

Consequently, the engine cannot search beyond 6 plies under standard time controls. Furthermore, the search core is single-threaded, preventing it from utilizing multi-core processor architectures.

\section{Threats to Validity}

\subsection{Internal Validity}
Internal validity threats concern variables that can affect timing and search statistics. In managed runtimes, background OS activity, JIT compiler warming, and collection sweeps introduce timing variations that can skew search latency benchmarks.

To mitigate these timing variations, all tournaments were executed on isolated hardware with background processes disabled. JIT warming was managed by executing initial dummy games to ensure V8 compiler optimization completed before recording metrics.

\subsection{External Validity}
External threats limit the generalizability of our results. The engine's rating scaling and pruning efficiencies were evaluated using a specific opening library (\texttt{positions.json}). Results could differ under rare or highly unbalanced opening configurations.

We addressed color and opening bias by enforcing paired color-swapped match runs (each opening is played twice, once as White and once as Black) across a randomized set of standard openings.

\subsection{Construct and Conclusion Validity}
Construct validity threats relate to our playing strength metrics. The relative Elo ratings reported in this work represent performance differences within a closed pool of Kronos variants rather than absolute playing strength. To calibrate these ratings, we benchmarked Kronos against depth-limited Stockfish 18 variants. However, because these matchups were restricted to 20 games per reference configuration due to execution latency, the absolute anchoring remains preliminary. Speculative projections for deeper Stockfish configurations should be treated as hypothetical rather than measured empirical data.

Finally, statistical conclusion validity is threatened by sample size constraints and high draw rates. The 100 to 400 game tournaments executed at depth 4 provide standard error margins of approximately $\pm 17$ to $\pm 34$ Elo. However, the high draw rates observed during self-play matches between similar variants (exceeding 92\% in some cases) significantly reduce the statistical information gain per game, widening the confidence intervals. While this is sufficient to confirm significance for major structural changes (like quiescence search or move ordering), smaller rating differentials require much larger game samples to resolve conclusively. We address this limitation by reporting the full LLR status and trinomial confidence intervals for all tournament outcomes, explicitly identifying cases where results remain inconclusive.


